\subsubsection{Reduction of General Systems}\label{subsec:general-system-reduction}

We now consider a general system of partial differential equations, for which we shall prescribe initial data on a noncharacteristic surface. We shall demonstrate how such a problem can be reduced to the form given by \eqref{eq:ck-intro-quasilinear-system}, \eqref{eq:ck-intro-zero-initial-data}. First, we reduce the initial surface to a plane.

A surface is usually described in one of three ways:
\begin{enumerate}
\item As a level set of a function: $\phi(x_1,\ldots,x_n)=\lambda$. Here it is assumed that $\nabla\phi\ne0$.
\item As a graph: $x_i=\psi(x_1,\ldots,x_{i-1},x_{i+1},\ldots,x_n)$.
\item By a parametrization, $\bm{x}=\bm{F}(\bm{y})$, where $\bm{y}\in\mathbb{R}^{n-1}$. Here it is assumed that the matrix $\partial F_i/\partial y_j$ has maximal rank.
\end{enumerate}

We say that a surface is of class $C^k$, $k\geq1$ (analytic) if, respectively, the functions $\phi$, $\psi$, $\bm{F}$ are of class $C^k$ (analytic). The implicit function theorem implies that locally all three definitions are equivalent. After possible renumbering of coordinates, we may therefore assume that an analytic surface is given in the form $x_n=\psi(x_1,\ldots,x_{n-1})$, where $\psi$ is analytic. We now set $y_i=x_i$ for $i\leq n-1$ and $y_n=x_n-\psi(x_1,\ldots,x_{n-1})$. In the new coordinates the initial surface is $y_n=0$. Henceforth we shall therefore assume that the initial surface is the plane $x_n=0$.

In reducing the equations, we shall sometimes have to differentiate equations with respect to $x_n$. We note that after such a differentiation, the plane $x_n=0$ remains a noncharacteristic surface of the new system. Moreover, if the differentiated equation is satisfied everywhere and the original equation is satisfied at $x_n=0$, then the original equation is satisfied everywhere. That is, whenever we differentiate an equation with respect to $x_n$, we shall view the original equation as a constraint which has to be satisfied by the initial data for the new system.

We now note that any system of PDEs can be made quasilinear by differentiating one or more of the equations with respect to $x_n$. Hence we can restrict our attention to quasilinear systems. Consider a system of the form
\begin{equation}\label{eq:general-weighted-quasilinear-system}
\sum_{j=1}^N\sum_{|\alpha|=s_i+t_j}a_{ij}^{\alpha}(\bm{p}_i)D^\alpha u_j+b_i(\bm{p}_i)=0,\qquad i=1,\ldots,N,
\end{equation}
where $s_i$ and $t_j$ are the weights assigned to the equations and dependent variables, respectively (cf. the discussion in Example~\ref{ex:weighted-laplace-reformulation}), and where the vector $\bm{p}_i$ consists of all the independent variables and all derivatives of the dependent variables $u_j$ of orders less than $s_i+t_j$. Since only the sums $s_i+t_j$ are relevant, we can assume that all the $s_i$ are non-positive and all the $t_j$ are non-negative. By differentiating the $i$th equation $-s_i$ times with respect to $x_n$, we can then make $s_i$ for the new system equal to zero. After doing so, a natural set of initial conditions is obtained by prescribing $\partial^l u^j/\partial x_n^l$ for $l=0,\ldots,t_j-1$. These initial conditions must be subjected to constraints arising from any equations which have been differentiated with respect to $x_n$.

At this point, we have arrived at a system of the form
\begin{equation}\label{eq:zero-row-weight-system}
\sum_{j=1}^N\sum_{|\alpha|=t_j}a_{ij}^{\alpha}(\bm{p})D^\alpha u_j+b_i(\bm{p})=0,\qquad i=1,\ldots,N,
\end{equation}
where $\bm{p}$ contains all the independent variables and all derivatives of all dependent variables $u_j$ of order less than $t_j$. The initial conditions are
\begin{equation}\label{eq:weighted-system-initial-derivatives}
\dfrac{\partial^l u_i}{\partial x_n^l}=f_{il}(x_1,\ldots,x_{n-1}),\qquad i=1,\ldots,N;\quad l=0,\ldots,t_i-1.
\end{equation}


Note that the vector $\bm{p}$ is determined by the initial data. If (for the given data) the surface $x_n=0$ is noncharacteristic, we can solve \eqref{eq:zero-row-weight-system} for the derivatives
\begin{equation}\label{eq:highest-normal-derivatives}
\dfrac{\partial^{t_i}u_i}{(\partial x_n)^{t_i}},
\end{equation}
and obtain a system of the form
\begin{equation}\label{eq:normal-solved-weighted-system}
\dfrac{\partial^{t_i}u_i}{(\partial x_n)^{t_i}}=\sum_{j=1}^N\sum_{\substack{|\alpha|=t_j\\\alpha\ne(0,\ldots,t_j)}}c_{ij}^{\alpha}(\bm{p})D^\alpha u_j+d_i(\bm{p}).
\end{equation}


We can now reduce this system to first-order by introducing all derivatives of $u_j$ which are of order less than $t_j$ as new dependent variables. We then add new equations accordingly, exploiting \eqref{eq:normal-solved-weighted-system} and the equality of mixed partial derivatives. For example if $t_1=2$, we would introduce the new variables $v_{1k}=\partial u_1/\partial x_k$ and add the equations
\begin{equation}\label{eq:reduction-first-order-derivative-relations}
\dfrac{\partial u_1}{\partial x_n}=v_{1n},\qquad\dfrac{\partial v_{1k}}{\partial x_n}=\dfrac{\partial v_{1n}}{\partial x_k},\qquad k=1,\ldots,n-1,
\end{equation}
and the equation for $\partial v_{1n}/\partial x_n$ would come from \eqref{eq:normal-solved-weighted-system}.

In this fashion, we arrive at an initial-value problem for a first-order system. The only differences from \eqref{eq:ck-intro-quasilinear-system}, \eqref{eq:ck-intro-zero-initial-data} which are left are that the initial data may not be zero and that the coefficients may depend on $x_n$. The former can be taken care of by a substitution. If the initial condition for $u_i$ is $u_i=f_i(x_1,\ldots,x_{n-1})$, we simply introduce $v_i=u_i-\allowbreak f_i(x_1,\ldots,x_{n-1})$ as new variables. Finally, we can remove the dependence of the coefficients on the independent variable $x_n$ by introducing an additional dependent variable $w$ which satisfies the equation
\begin{equation}\label{eq:reduction-extra-independent-coordinate}
\dfrac{\partial w}{\partial x_n}=1,\qquad w=0\quad\text{on }x_n=0.
\end{equation}
In all coefficients which depend on $x_n$, we then replace $x_n$ by $w$.

\begin{Example}\label{ex:scalar-general-reduction}Let us consider the scalar equation
\begin{equation}\label{eq:scalar-nonlinear-reduction-origin}
\phi(u_{xx},u_{xy},u_{yy},u_x,u_y,u,x,y)=0
\end{equation}
with initial conditions
\begin{equation}\label{eq:scalar-nonlinear-reduction-initial}
u(x,0)=f(x),\qquad u_y(x,0)=g(x).
\end{equation}
We assume that $\phi$, $f$ and $g$ are analytic functions of their arguments and that, at least in a neighborhood of the origin, the plane $y=0$ is noncharacteristic, i.e., the equation
\begin{equation}\label{eq:scalar-nonlinear-initial-constraint}
\phi(f''(x),g'(x),z,f'(x),g(x),f(x),x,0)=0
\end{equation}
has a solution $z(x)$, and
\begin{equation}\label{eq:scalar-nonlinear-noncharacteristic}
\phi_{,3}(f''(x),g'(x),z(x),f'(x),g(x),f(x),x,0)\ne0.
\end{equation}
Here and in the following, $\phi_{,i}$ denotes the derivative of $\phi$ with respect to the $i$th argument.

We first make the equation quasilinear by differentiating with respect to $y$. This leads to the problem
\begin{equation}\label{eq:differentiated-scalar-nonlinear-equation}
\begin{aligned}
0&=\phi_{,1}(\cdots)u_{xxy}+\phi_{,2}(\cdots)u_{xyy}+\phi_{,3}(\cdots)u_{yyy}\\
&\quad+\phi_{,4}(\cdots)u_{xy}+\phi_{,5}(\cdots)u_{yy}\\
&\quad+\phi_{,6}(\cdots)u_y+\phi_{,8}(\cdots),
\end{aligned}
\end{equation}
(where the dots $(\cdots)$ denote the arguments $(u_{xx},u_{xy},u_{yy},u_x,u_y,u,x,y)$) and the initial conditions are
\begin{equation}\label{eq:differentiated-scalar-initial-data}
u(x,0)=f(x),\quad u_y(x,0)=g(x),\quad u_{yy}(x,0)=z(x),
\end{equation}
with $z(x)$ as above.

We now introduce new variables $p=u_x$, $q=u_y$, $a=u_{xx}$, $b=u_{xy}$, $c=u_{yy}$, $y=w$ and obtain the system
\begin{equation}\label{eq:first-order-scalar-nonlinear-system}
\begin{aligned}
u_y&=q,\quad p_y=b,\quad q_y=c,\quad a_y=b_x,\quad b_y=c_x,\\
c_y&=-\dfrac1{\phi_{,3}(\cdots)}
\bigl[\phi_{,1}(\cdots)b_x+\phi_{,2}(\cdots)c_x\\
&\qquad+\phi_{,4}(\cdots)b+\phi_{,5}(\cdots)c+\phi_{,6}(\cdots)q+\phi_{,8}(\cdots)\bigr],\\
w_y&=1.
\end{aligned}
\end{equation}
Here the dots $(\cdots)$ indicate the argument $(a,b,c,p,q,u,x,w)$. The initial conditions are
\begin{equation}\label{eq:first-order-scalar-nonlinear-initial-data}
\begin{gathered}
u(x,0)=f(x),\quad p(x,0)=f'(x),\quad q(x,0)=g(x),\quad a(x,0)=f''(x),\\
b(x,0)=g'(x),\quad c(x,0)=z(x),\quad w(x,0)=0.
\end{gathered}
\end{equation}
\end{Example}

\begin{Example}\label{ex:stokes-general-reduction}As an example for a system, consider the Stokes system in two space dimensions,
\begin{equation}\label{eq:stokes-two-dimensional-reduction-origin}
\begin{aligned}
u_{xx}+u_{yy}-p_x&=0,\\v_{xx}+v_{yy}-p_y&=0,\\u_x+v_y&=0,
\end{aligned}
\end{equation}
with initial conditions on $y=0$. With the understanding that the ordering of variables is $(u,v,p)$, we can assign the weights $s_1=s_2=0$, $s_3=-1$, $t_1=t_2=2$, $t_3=1$. According to the above algorithm, we must differentiate the last equation in \eqref{eq:stokes-two-dimensional-reduction-origin} with respect to $y$, leading to
\begin{equation}\label{eq:stokes-differentiated-divergence}
u_{xy}+v_{yy}=0.
\end{equation}
The system can then be restated as
\begin{equation}\label{eq:stokes-normal-solved-system}
\begin{aligned}
u_{yy}&=-u_{xx}+p_x,\\v_{yy}&=-u_{xy},\\p_y&=v_{xx}-u_{xy}.
\end{aligned}
\end{equation}


Appropriate initial conditions are
\begin{equation}\label{eq:stokes-reduction-initial-data}
\begin{gathered}
u(x,0)=f_1(x),\quad u_y(x,0)=f_2(x),\\
v(x,0)=g_1(x),\quad v_y(x,0)=g_2(x),\\
p(x,0)=h(x);
\end{gathered}
\end{equation}
however, in order to obtain a solution of the original system, these initial data must satisfy the constraint given by the last equation of \eqref{eq:stokes-two-dimensional-reduction-origin}, namely,
\begin{equation}\label{eq:stokes-reduction-initial-constraint}
f_1'(x)+g_2(x)=0.
\end{equation}
To obtain a first-order system, we finally set $u_x=a$, $u_y=b$, $v_x=c$, $v_y=d$. We obtain
\begin{equation}\label{eq:stokes-first-order-reduction}
\begin{gathered}
u_y=b,\quad a_y=b_x,\quad b_y=-a_x+p_x,\quad v_y=d,\\
c_y=d_x,\quad d_y=-b_x,\quad p_y=c_x-b_x.
\end{gathered}
\end{equation}
\end{Example}
